\documentclass[../../../include/open-logic-section]{subfiles}

\begin{document}

\olfileid{sth}{z}{union}
\olsection{اتحاد}

\olref[sth][z][sep]{prop:intersectionsexist} تقاطعګانې راکړې.
خو که غواړو چې خپل‌سري اتحادونه شتون ولري، نو يو بل بديهي
اصل هم ټاکل پکار دي:

\begin{axiom}[اتحاد] د هر سټ $A$ دپاره سټ $\bigcup A =
\Setabs{x}{(\exists b \in A) x \in b}$ شتون لري.
\[
	\forall A \exists U \forall x(x \in U \liff (\exists b \in A)x \in b)
\]
\end{axiom}

د تجمعي-تکراري تصور له مخې دا بديهي اصل هم توجيه کېږي.
$A$ يو سټ وګرځوئ؛ نو $A$ په يوه پړاو $S$ کښې جوړېږي
(د \stageshier له مخې). د $A$ هر غړے له $S$ \emph{نه مخکې}
جوړ شوے دے (د \stagesacc له مخې)؛ نو په همدې ډول استدلال
کولو سره، د $A$ د هر غړي هر غړے له $S$ نه مخکې جوړ شوے دے.
په دې توګه دا \emph{ټول} سټونه له $S$ نه مخکې موجود دي،
څو په $S$ کښې په يوه سټ کښې راټول شي. او همدغه سټ $\bigcup A$ دے.

\end{document}
